% =============================================================================
% 2 Background and Motivation. Written 2026-10-20 from the approved outline
% (paper_versions/sec0-2-outline.md). Prefill appears only as the other half of
% P/D disaggregation. The "when worthwhile" figure uses the A100 M13 sweep
% only (paper/data/overhead/15_clean_cf/CF_c128_a.json) and says nothing about
% how c depends on context: that question is open (M22, analysed with Ben).
% =============================================================================

\section{Background and Motivation}\label{sec:background}

% Restructured 2026-10-20 (Ben's review): six short subsections merged into
% three, following the argument -- the decode node is HBM-bound; host memory
% over a fast link is a second tier that existing systems use one path at a
% time; when riders pay.
\subsection{The Decode Node Is Bound by HBM Capacity}

A transformer layer's attention compares each new token's query with the keys
of every earlier token and mixes their values. Serving engines keep those keys
and values, the KV cache, so that a decode step computes them only for the new
token. A model with $L$ layers, $h_{\mathrm{kv}}$ KV heads of dimension $d$
and $s$ bytes per element stores $\btok = 2L\,h_{\mathrm{kv}}\,d\,s$ bytes per
token: 320\,KiB for Qwen2.5-72B with 16-bit KV. A decode step over a batch of
requests with contexts $S_i$ reads the weights $W$ once and
$\sum_i S_i\btok$ bytes of KV. The weights are shared by the batch; the KV is
not.
% src: Qwen2.5-72B config: 80 layers, 8 KV heads, head dim 128 ->
%      2*80*8*128*2 = 327,680 B (paper/model/overhead_model.py Q72).

Prefill and decode stress different resources, so production systems place
them on separate nodes and hand each request's KV from one to the
other~\cite{distserve,splitwise,mooncake}. This paper is about the decode
node. Its KV arrives over the network once per request, and where it lands,
HBM or host memory, is the decode node's choice.

% Section 2.3: the compute a decode step can use, at the largest batch that fits,
% as a roofline bound. At max batch B(S) = M_kv/(S b_tok) a step reads W + M_kv
% bytes and does 2 P B FLOPs, so utilisation <= (2 P B / (W + M_kv)) / ridge.
% Analytic, from paper/model/overhead_model.py (Q72: W 76.0 GB, P 72.7e9,
% b_tok 327,680 B; Q32: W 65.53 GB, P 32.8e9, b_tok 262,144 B; 4 GB reserve).
% GH200: fp8 dense 1979 TFLOP/s / 4 TB/s = ridge 495. A100: fp16 312 TFLOP/s /
% 1.935 TB/s = ridge 161. Computed 2026-10-20 (outline sec0-2, 2.3).
\begin{figure}[t]
\centering
\begin{tikzpicture}
\begin{loglogaxis}[
  width=0.95\columnwidth, height=4.6cm,
  xlabel={context $S$ (tokens)}, ylabel={compute usable (\%, bound)},
  xmin=800, xmax=170000, ymin=0.08, ymax=40,
  xtick={1024,4096,16384,65536}, xticklabels={1k,4k,16k,64k},
  ytick={0.1,1,10}, yticklabels={0.1,1,10},
  tick label style={font=\sffamily\scriptsize},
  label style={font=\sffamily\scriptsize},
  legend style={font=\sffamily\scriptsize, draw=none, fill=white,
                fill opacity=0.85, text opacity=1, at={(0.98,0.97)},
                anchor=north east, inner sep=1pt, row sep=-1pt},
  legend cell align=left,
  axis line style={black!60}, tick style={black!60},
  ymajorgrids, grid style={black!10},
  every axis plot/.append style={line width=0.9pt},
]
\addplot[hbmD, mark=*, mark size=1.4pt] table[x=S, y=u] {
S       u
1024    15.23
2048    7.62
4096    3.81
8192    1.90
16384   0.95
32768   0.48
65536   0.24
131072  0.12
};
\addlegendentry{GH200, Qwen2.5-72B fp8}
\addplot[ddrD, mark=square*, mark size=1.3pt, densely dashed] table[x=S, y=u] {
S       u
1024    20.88
2048    10.44
4096    5.22
8192    2.61
16384   1.30
32768   0.65
65536   0.33
131072  0.16
};
\addlegendentry{A100, Qwen2.5-32B fp16}
\node[font=\sffamily\scriptsize, text=hbmD, anchor=north east]
  at (axis cs:16384,0.80) {3 requests};
\node[font=\sffamily\scriptsize, text=hbmD, anchor=south west]
  at (axis cs:1100,15.5) {48 requests};
\end{loglogaxis}
\end{tikzpicture}
\caption{\textbf{Decode leaves compute idle, and more so as context grows.}
Roofline bound on the share of the GPU's compute a decode step can use at the
largest batch that fits in the HBM left after the weights. Analytic.}
\label{fig:util}
\end{figure}


On that node, because the weights are read once per step, the work per byte read grows with
the batch, and so does the share of the GPU's compute a step can use. The
batch, in turn, is what fits in the HBM the weights leave free: at most
$\Mkv/(S\btok)$ requests of context $S$. \Cref{fig:util} combines the two into
a roofline bound. For Qwen2.5-72B in fp8 on a GH200, a full batch at 1k tokens
can use at most 15\% of the GPU's compute; at 16k tokens it is three requests
and under 1\%. The bound falls as $1/S$. Measurements of large-batch decode
agree that compute, not memory, is what goes unused~\cite{mindgap}.
% src: figs/fig-util.tex data (paper/model/overhead_model.py constants; fp8
%      dense peak 1979 TFLOP/s, HBM 4 TB/s, ridge 495 FLOP/B).

\subsection{A Second Memory Tier}

The second tier is reached over a link: PCIe on most servers, NVLink-C2C on a
GH200 (450\,GB/s per direction~\cite{gh200page,directkv}, 419\,GB/s measured
from host to device~\cite{pie}), and a PCIe-class link to a CXL memory
expander. Whatever the tier, the bandwidth available to a GPU is that of the
slower of the GPU's own host link and the tier's link; we call it $\Xlink$.
% src: directkv s2: "CPU-GPU transfers (900 GB/s, 450 GB/s per direction)";
%      pie s2: "It achieves a data transfer rate of 419 GB/s from CPU to GPU".
% TODO-verify: CXL 3.x and 4.0 per-x16 bandwidth (derived 128 / 256 GB/s per
%      direction from 64 / 128 GT/s) awaits the CXL Consortium specification.

% Section 2.5 (Ben's diagram request): which tier serves which request, per
% system. Six requests R1..R6 in every panel. Drawn from the systems' own text:
%   (a) Pie s3.1 "We relocate m layers of the KV cache to the CPU, leaving n-m
%       layers on the GPU"; s3.3 swap in and swap out through GPU buffers.
%   (b) DirectKV s6 all K,V "in CPU memory"; abstract "GPU kernels to directly
%       access CPU-resident KV cache"; s1 streamed through SMEM.
%   (c) HeteroKV, s4: residents in HBM, riders in host DDR streamed into a
%       k-slot window, one forward pass; the link budget is X T0 per step.
% DAK, which also splits by request (concurrent work), is discussed in s8, not
% drawn. Colours from figstyle: HBM teal, DDR orange, H2D blue, D2H purple.
\begin{figure*}[t]
\centering
\begin{tikzpicture}[
  req/.style={draw=#1D, fill=#1F, line width=0.5pt, minimum width=5.6mm,
              minimum height=4.2mm, inner sep=0pt, font=\sffamily\tiny},
  half/.style={draw=#1D, fill=#1F, line width=0.5pt, minimum width=5.6mm,
               minimum height=2.1mm, inner sep=0pt},
  wts/.style={draw=black!45, fill=black!10, minimum height=4.2mm,
              inner sep=1pt, font=\sffamily\tiny},
  fwd/.style={cmp, minimum height=4.6mm, inner sep=1pt, font=\sffamily\scriptsize},
  tier/.style={font=\sffamily\scriptsize\bfseries, anchor=east},
  note/.style={font=\sffamily\scriptsize, text=inkS, align=center},
]
% ---------------------------------------------------------------- layout
% panel x offsets 0, 6.0, 12.0 cm; rows: compute 3.2, HBM 2.2, link 1.2, DDR 0.2
\foreach \px/\ttl in {0/{(a) Layer swap via HBM buffer},
                      6.0/{(b) Zero-copy (DirectKV)},
                      12.0/{(c) Split by request (\sysname{})}} {
  \begin{scope}[xshift=\px cm]
    \fill[bgA, rounded corners=3pt] (-0.15,-0.25) rectangle (5.55,3.75);
    \node[hd, anchor=south] at (2.7,3.78) {\ttl};
    \node[tier] at (0.9,2.2) {HBM};
    \node[tier] at (0.9,0.2) {DDR};
    \node[tier, font=\sffamily\scriptsize] at (0.9,1.2) {link};
  \end{scope}
}
% ------------------------------------------------------------------ (a) Pie
\begin{scope}[xshift=0cm]
  \node[fwd, minimum width=4.3cm] (fa) at (3.15,3.2) {one forward pass, R1--R6};
  \node[wts, minimum width=1.0cm] at (1.55,2.2) {weights};
  \foreach \i in {1,...,6} {
    \node[half=hbm] (ah\i) at ({1.95+0.5*\i},2.31) {};
    \node[half=ddr] (ad\i) at ({1.95+0.5*\i},0.09) {};
    \node[font=\sffamily\tiny] at ({1.95+0.5*\i},-0.1) {R\i};
  }
  \node[plain, minimum width=1.1cm, minimum height=2.1mm, inner sep=0pt,
        font=\sffamily\tiny] (abuf) at (3.7,1.95) {};
  \node[font=\sffamily\tiny, anchor=east] at (3.13,1.95) {buffer};
  \draw[h2d] (3.45,0.25) -- (3.45,1.84);
  \draw[d2h] (3.95,1.84) -- (3.95,0.25);
  \node[font=\sffamily\tiny, text=d2hC, anchor=west] at (4.0,1.05) {swap out};
  \draw[flow] (3.15,2.44) -- (3.15,2.96);
  \node[note] at (2.7,-0.05) {};
\end{scope}
% ------------------------------------------------------------- (b) DirectKV
\begin{scope}[xshift=6.0cm]
  \node[fwd, minimum width=4.3cm] (fb) at (3.15,3.2) {one forward pass, R1--R6};
  \node[wts, minimum width=1.0cm] at (1.55,2.2) {weights};
  \node[draw=black!40, densely dashed, minimum width=2.9cm, minimum height=4.2mm,
        font=\sffamily\tiny, text=inkS] at (3.65,2.2) {no KV in HBM};
  \foreach \i in {1,...,6} {
    \node[req=ddr] (bd\i) at ({1.95+0.5*\i},0.2) {R\i};
  }
  \draw[h2d, line width=1.3pt] (4.9,0.45) .. controls (5.35,1.4) and (5.35,2.6)
    .. (4.95,2.95);
  \node[font=\sffamily\tiny, text=h2dC, anchor=east, align=right]
    at (5.0,1.5) {straight into\\SMEM};
\end{scope}
% -------------------------------------------------------------- (c) HeteroKV
\begin{scope}[xshift=12.0cm]
  \node[fwd, minimum width=4.3cm] (fc) at (3.15,3.2) {one forward pass, R1--R5};
  \node[wts, minimum width=1.0cm] at (1.55,2.2) {weights};
  \foreach \i in {1,2,3} {
    \node[req=hbm] (ch\i) at ({1.95+0.5*\i},2.2) {R\i};
  }
  \node[draw=hbmD, fill=white, line width=0.5pt, minimum width=1.05cm,
        minimum height=4.2mm, inner sep=0pt, font=\sffamily\tiny]
    (cwin) at (4.55,2.2) {$k$ slots};
  \foreach \i in {4,5} {
    \node[req=ddr] (cd\i) at ({2.95+0.5*(\i-4)+0.75},0.2) {R\i};
  }
  \node[draw=black!40, densely dashed, fill=white, minimum width=5.6mm,
        minimum height=4.2mm, inner sep=0pt, font=\sffamily\tiny, text=inkS]
    (cd6) at (5.05,0.2) {R6};
  \draw[h2d] (4.0,0.45) -- (4.3,1.98);
  \draw[h2d] (4.45,0.45) -- (4.6,1.98);
  \draw[decorate, decoration={brace, amplitude=2.5pt, mirror}, black!60]
    (3.4,0.75) -- (3.4,1.65);
  \node[font=\sffamily\tiny, anchor=east, align=right] at (3.32,1.2)
    {at most $\Xlink T_0$\\bytes per step};
  \draw[flow] (2.95,2.44) -- (2.95,2.96);
  \draw[flow] (4.55,2.44) -- (4.55,2.96);
  \node[font=\sffamily\tiny, text=inkS, anchor=north] at (5.05,-0.02) {waits};
\end{scope}
\end{tikzpicture}
\caption{\textbf{One path per byte, or both.} (a) Layer swapping~\cite{pie} keeps a
share of every request's layers, sized online, in host memory and swaps them through an HBM buffer, in
and back out. (b) DirectKV keeps all KV in host memory and streams it into the
kernel's shared memory; HBM holds no KV. (c) \sysname{} splits the batch by
request: residents (R1--R3) read HBM, riders (R4, R5) stream from DDR into a
$k$-slot window within the link's per-step budget, and both run in one forward
pass. R6 waits.}
\label{fig:tiers}
\end{figure*}


Existing systems that use host memory for KV send each byte through one path
(\cref{fig:tiers}). Layer swapping keeps a share of every request's layers in
host memory, sized online to what the link can hide, and swaps them through
an HBM buffer~\cite{pie}. Every request crosses the link every step, and each layer is
swapped back out after use. DirectKV keeps all KV in host memory and streams it into the
attention kernel's shared memory~\cite{directkv}. Its HBM holds weights and
no KV. In both, a byte of KV is read either at HBM bandwidth or at link
bandwidth; the two are never used for different requests in the same step.
% src: pie s3.1 (per-layer relocation, quoted in s1); s3.3: "Each layer takes
%      the equivalent of two layer computation times to swap in ... and another
%      two layer computation times to swap out". directkv s6 (all KV in CPU
%      memory, quoted in s1); abstract: "enabling GPU kernels to directly
%      access CPU-resident KV cache".
%      pie s1 adaptive expansion: "It starts with zero CPU memory and then
%      gradually increases this memory as long as ... the GPU-CPU interconnect
%      is not saturated".

% Section 2.5: the batch riders could add, r*/B = T0 X / M_kv, against the share
% of HBM the weights take, one curve per link bandwidth X. Reference GPU: 96 GB
% HBM at 4 TB/s, eta 0.75, 4 GB activation reserve; T0 = W/(eta BW),
% M_kv = 96 - W - 4. Computed 2026-10-20 (outline sec0-2, 2.5). 25.1 GB/s is
% the A100 PCIe4 copy rate we measure; 419 GB/s is Pie's measured GH200 C2C.
% 128 and 256 GB/s are plain bandwidths. TODO-verify: they correspond to CXL
% 3.x and 4.0 x16 links only once the CXL Consortium specification is cited.
% Target: Qwen2.5-72B fp8 on a GH200-96, w = 76/96 = 0.79, r*/B = 0.663.
\begin{figure}[t]
\centering
\begin{tikzpicture}
\begin{semilogyaxis}[
  width=0.86\columnwidth, height=5.0cm,
  xlabel={share of HBM taken by the weights},
  ylabel={riders per resident, $r_0/B$},
  xmin=0.28, xmax=0.92, ymin=0.003, ymax=4,
  xtick={0.3,0.4,0.5,0.6,0.7,0.8,0.9},
  xticklabels={30\%,40\%,50\%,60\%,70\%,80\%,90\%},
  ytick={0.01,0.1,1}, yticklabels={0.01,0.1,1},
  tick label style={font=\sffamily\scriptsize},
  label style={font=\sffamily\scriptsize},
  axis line style={black!60}, tick style={black!60},
  ymajorgrids, grid style={black!10},
  every axis plot/.append style={line width=0.9pt},
  clip=false,
]
\pgfplotstableread{
w     p4     p5     x128   x256   c2c
0.30  0.004  0.010  0.019  0.039  0.064
0.40  0.006  0.015  0.031  0.061  0.100
0.50  0.009  0.023  0.047  0.093  0.152
0.60  0.014  0.036  0.071  0.143  0.234
0.70  0.023  0.058  0.116  0.231  0.378
0.75  0.030  0.077  0.154  0.307  0.503
0.79  0.039  0.100  0.200  0.400  0.655
0.83  0.054  0.138  0.276  0.552  0.903
0.87  0.082  0.210  0.420  0.840  1.376
0.90  0.129  0.329  0.658  1.317  2.155
}\potential
\addplot[black!45]                 table[x=w, y=p4]   {\potential};
\addplot[black!70, densely dashed] table[x=w, y=p5]   {\potential};
\addplot[h2dC!60]                  table[x=w, y=x128] {\potential};
\addplot[h2dC, densely dashed]     table[x=w, y=x256] {\potential};
\addplot[hbmD, line width=1.2pt]   table[x=w, y=c2c]  {\potential};
\node[font=\sffamily\scriptsize, text=black!55, anchor=west] at (axis cs:0.905,0.129) {25\,GB/s};
\node[font=\sffamily\scriptsize, text=black!75, anchor=west] at (axis cs:0.905,0.329) {64};
\node[font=\sffamily\scriptsize, text=h2dC!70, anchor=west]  at (axis cs:0.905,0.62) {128};
\node[font=\sffamily\scriptsize, text=h2dC, anchor=west]     at (axis cs:0.905,1.2) {256};
\node[font=\sffamily\scriptsize, text=hbmD, anchor=west]     at (axis cs:0.905,2.45) {419};
\addplot[only marks, mark=*, mark size=2.2pt, draw=black, fill=ddrD]
  coordinates {(0.7917,0.663)};
\node[font=\sffamily\scriptsize, anchor=east, align=right] (tg)
  at (axis cs:0.70,1.7) {GH200, Qwen2.5-72B fp8};
\draw[black!60, line width=0.5pt] (tg.east) -- (axis cs:0.785,0.72);
\end{semilogyaxis}
\end{tikzpicture}
\caption{\textbf{What the link could add.} Requests the link can carry per step
relative to those HBM holds, $r_0/B = T_0\Xlink/\Mkv$, for a GPU with
96\,GB of HBM at 4\,TB/s, one curve per link bandwidth $\Xlink$ (GB/s; 25 and
419 are measured rates). An upper bound, before the costs $F$ and $c$ of
\S\ref{sec:model}.}
\label{fig:potential}
\end{figure}


\Cref{fig:tiers}c is the alternative this paper builds. The batch is split by
request. Residents fill HBM as before. Riders keep their KV in host memory and
stream it in one layer ahead, as many as the link can carry within a step:
$\Xlink T_0$ bytes, for a step time $T_0$. For a workload of one context $S$,
HBM holds $B = \Mkv/(S\btok)$ residents and the link carries
$r_0 = T_0\Xlink/(S\btok)$ riders, so riders add
$r_0/B = T_0\Xlink/\Mkv$ to the batch. \Cref{fig:potential} plots that
ratio. It is large only when the link is fast and the weights take most of
HBM, since both lengthen $T_0$ and shrink $\Mkv$. On a GH200 serving
Qwen2.5-72B in fp8, the link could carry two thirds as many requests again as
HBM holds. This is an upper bound: it ignores what a rider costs.
% src: figs/fig-potential.tex data: w = 76/96, T0 = 76/(0.75*4000) = 25.3 ms,
%      M_kv = 16 GB, X = 419 GB/s -> r*/B = 0.663.

\subsection{When Riders Pay}

% Section 2.6: when riders pay, on one measured A100 sweep. Data: M13,
% CF_c128_a (job 2181931, paper/data/overhead/15_clean_cf/CF_c128_a.json):
% Qwen2.5-32B fp16, 16 residents at context 128, riders 1..8, T0 = 47.74 ms.
% Fit from figs/data/fits.txt (abl_real_a): F = 6.050, c = 0.508.
% Token worth T0/B = 2.98 ms per rider; r_min = 3; knee = T0 / link-per-rider
% = 47.74 / 1.696 = 28.1 riders (link at 25.1 GB/s). Past the knee the curve is
% the model F + c r + (r link - T0), NOT measured, and drawn dashed.
% NO M22 data, and no claim about how c depends on context (Ben, 2026-10-20).
\begin{figure}[t]
\centering
\begin{tikzpicture}
\begin{axis}[
  width=0.95\columnwidth, height=5.0cm,
  xlabel={riders $r$}, ylabel={added step time (ms)},
  xmin=0, xmax=32, ymin=0, ymax=60,
  xtick={0,4,8,12,16,20,24,28,32},
  tick label style={font=\sffamily\scriptsize},
  /pgf/number format/assume math mode=true,
  label style={font=\sffamily\scriptsize},
  axis line style={black!60}, tick style={black!60},
  ymajorgrids, grid style={black!10},
  every axis plot/.append style={line width=0.9pt},
  legend style={font=\sffamily\scriptsize, draw=none, fill=none,
                at={(0.5,-0.30)}, anchor=north, inner sep=1pt,
                legend columns=3, column sep=4pt},
  legend cell align=left,
]
% riders pay where the overhead lies below the tokens' worth
\addplot[name path=worth, draw=none, forget plot, domain=2.44:28.1, samples=2]
  {2.984*x};
\addplot[name path=over, draw=none, forget plot, domain=2.44:28.1, samples=2]
  {6.050+0.508*x};
\addplot[cmpF, forget plot] fill between[of=worth and over];
\addplot[cmpD, densely dashed, domain=0:20, samples=2] {2.984*x};
\addlegendentry{tokens' worth, $r\,T_0/B$}
\addplot[ddrD, domain=0:28.1, samples=2] {6.050+0.508*x};
\addlegendentry{$F + c\,r$, fitted}
\addplot[ddrD, mark=*, mark size=1.5pt, only marks]
  table[x=r, y=dT, col sep=comma] {figs/data/abl_real_a.csv};
\addlegendentry{measured (A100)}
\addplot[ddrD, densely dotted, domain=28.1:32, samples=2, forget plot]
  {6.050+0.508*x + (1.696*x-47.74)};
% knee and break-even
\draw[black!50, densely dashed, line width=0.6pt]
  (axis cs:28.1,0) -- (axis cs:28.1,44);
\node[font=\sffamily\scriptsize, text=black!70, anchor=south, align=center]
  at (axis cs:28.1,44) {knee\\$r^\star{=}28$};
\draw[black!50, line width=0.6pt] (axis cs:3,0) -- (axis cs:3,3.4);
\node[font=\sffamily\scriptsize, text=black!70, anchor=south west, inner sep=1pt]
  at (axis cs:3.2,0.3) {$r_{\min}{=}3$};
\node[font=\sffamily\scriptsize, text=cmpD, anchor=center]
  at (axis cs:16,25) {riders pay};
\node[font=\sffamily\scriptsize, text=ddrD, anchor=south west, align=left]
  at (axis cs:12.5,0.4) {intercept $F{=}6.05$\\slope $c{=}0.51$};
\end{axis}
\end{tikzpicture}
\caption{\textbf{When riders pay}, on one A100 sweep (Qwen2.5-32B, 16
residents, $T_0 = 47.7$\,ms). Each rider adds $c$ and the first adds $F$; each
is worth $T_0/B$ of step time. Riders pay in the shaded region, from
$\rmin$ to the knee, where their transfer stops hiding behind the step; the dotted line past the knee is the model,
not a measurement.}
\label{fig:worth}
\end{figure}


A rider joins the batch as one more row of every matrix multiply; only its
attention reads a different place. It is not free. The first rider adds a
fixed cost $F$ to the step, and each rider adds a further cost $c$. A batch of
$B$ residents produces $B$ tokens per step, so one token is worth $T_0/B$ of
step time, and riders pay when they add more tokens than step time:
$r\,(T_0/B - c) > F$. Two conditions follow. A rider must cost less than a
token is worth, $c < T_0/B$. And there must be enough riders to repay $F$.
Riders can be added only while their KV crosses the link within one step;
beyond that point, the \emph{knee}, each rider's transfer adds to the step.

\Cref{fig:worth} shows the three on one measured A100 run. Below the dashed
line, the riders' tokens are worth more than the time they add; three riders
are enough to repay $F$ here, and the knee lies far to the right.
\S\ref{sec:model} turns these conditions into a model, and
\S\ref{sec:ablation} measures what $F$ and $c$ are made of.
% src: CF_c128_a (job 2181931): T0 = 47.74 ms, B = 16, F = 6.050, c = 0.508
%      (figs/data/fits.txt abl_real_a); T0/B = 2.98; r_min = floor(6.05 /
%      (2.98 - 0.508)) + 1 = 3; link 1.696 ms per rider -> knee 47.74/1.696 =
%      28.1 riders.
